\chapter{Userspace Drivers}

In CharlotteOS, device drivers are userspace processes. They receive only the
resources they need, they communicate through endpoints, and their failures are
contained. This chapter explains the userspace driver model.

\section{The problem with kernel drivers}

In a monolithic kernel, a driver has access to everything: all physical memory,
all interrupt vectors, all kernel data structures. A bug in a driver --- a buffer
overflow, a use-after-free, a race condition --- can corrupt the kernel and crash
the system. This is not merely a theoretical concern: driver bugs are the dominant cause
of operating system crashes in most operating systems.

The microkernel solution is to move drivers to userspace. But simply moving
them is not enough: a userspace driver that can still map arbitrary physical
memory or receive arbitrary interrupts is still dangerous. The driver must be
\textbf{granted} only what it needs, and the grants must be mediated by the
kernel.

\section{The DriverGrant}

When a driver domain is created, the supervisor delivers a \texttt{DriverGrant}
through the config-page contract:

\begin{verbatim}
struct DriverGrant {
    device: DeviceCap,        // identity and lifecycle of the device
    mmio: Vec<MmioRegionCap>, // register windows, page-granular
    interrupts: Vec<InterruptCap>, // IRQ sources for this device
    dma_domain: Option<DmaDomainCap>, // DMA translation, if IOMMU present
    bootstrap_endpoint: EndpointCap,  // endpoint on which to serve clients
}
\end{verbatim}

The driver receives \textbf{exactly} the MMIO register windows and interrupt
vectors it needs. It cannot:
\begin{itemize}
    \item Map arbitrary physical memory.
    \item Receive interrupts not assigned to its device.
    \item Access other devices' MMIO regions.
    \item Perform DMA without explicit pinning into its DMA domain.
\end{itemize}

\section{MMIO delegation}

An \texttt{MmioRegion} capability represents a page-granular device register
window. The driver maps it into its address space as Device-nGnRnE memory
(user-accessible, execute-never, strongly ordered on AArch64; uncacheable on
x86-64).

The mapping is a kernel-mediated operation: the driver calls
\texttt{mmio\_map(region\_cap, virtual\_address)}. The kernel validates the
capability, creates the page-table entries with the correct memory attributes,
and returns. The driver then performs \textbf{direct EL0 MMIO} --- load and
store instructions to the mapped virtual addresses, with no kernel involvement
on the fast path.

This is the key performance property: the common case (reading a device register,
writing to a FIFO) is one CPU instruction. There is no syscall overhead per
MMIO access.

\section{Interrupt delivery}

An \texttt{Interrupt} capability represents a routed interrupt source. The driver
binds it to its completion queue:

\begin{verbatim}
interrupt_bind(int_cap, cq_id);
\end{verbatim}

When the interrupt fires:

\begin{enumerate}
    \item The IRQ dispatcher reads the per-interrupt route table atomically.
    \item It masks the source at the interrupt controller (GICv3 on AArch64,
        x2APIC on x86-64).
    \item It atomically bumps a coalescing counter.
    \item It pushes a wake onto a deferred queue. The actual
        \texttt{completion::wake} (which takes locks) is performed from
        thread context by \texttt{yield\_lp} and the LP idle loop.
    \item The driver's shard wakes.
    \item The driver reads the device's status register (direct EL0 MMIO) to
        determine the interrupt cause.
    \item The driver processes the interrupt and unmasks the source.
\end{enumerate}

This is a lock-free fast path: the interrupt context does very little work (atomic
counter bump, deferred-wake enqueue), deferring the rest to thread context.

\section{The reference UART driver}

The PL011 UART driver demonstrates the full lifecycle:

\begin{itemize}
    \item \textbf{Direct MMIO writes to the transmit FIFO.} No syscalls.
    \item \textbf{Interrupt-driven deferred reads.} The driver retains a reply
        token for a pending read. When the RX interrupt fires, the driver reads
        the receive register and completes the deferred reply.
    \item \textbf{Cooperative shutdown.} On graceful exit, the driver unmap its
        MMIO region and closes its interrupt.
    \item \textbf{Crash recovery.} If the driver terminates without releasing
        its device caps (panic, infinite loop, killed by supervisor), the
        teardown path:
        \begin{enumerate}
            \item Reclaims the MMIO region (unmaps it, returns to the kernel).
            \item Unroutes the interrupt.
            \item Reconciles outstanding operations: any pending reply tokens
                are invalidated; any borrowed memory is returned to lenders.
            \item Resets the device hardware.
        \end{enumerate}
        A restarted instance with fresh grants resumes under a bumped
        generation.
\end{itemize}

\section{DMA and the IOMMU}

For DMA-capable devices (network cards, storage controllers), a userspace
driver needs more: the ability to tell the device ``read from physical address
X'' or ``write to physical address Y.'' Without an IOMMU, this is a safety
gap: a malicious or buggy driver could program the device to DMA into kernel
memory or another domain's pages.

The architecture anticipates two levels of isolation:

\begin{itemize}
    \item \textbf{Without IOMMU:} the driver is trusted with DMA. It still
        receives only the MMIO and interrupt it needs, and it still runs in
        a separate address space, so a CPU-side bug cannot read kernel memory.
        But a DMA programming bug can violate memory isolation.
    \item \textbf{With IOMMU/SMMU:} the DMA domain capability binds the
        device to a specific IOMMU translation table. The driver can only
        DMA into pages explicitly pinned into its DMA domain. The kernel
        validates every pin/unpin through the capability. The IOMMU hardware
        enforces the rest --- even a malicious driver cannot DMA outside its
        domain.
\end{itemize}

The driver API is the same in both cases: \texttt{dma\_pin(memory\_object)} and
\texttt{dma\_unpin(memory\_object)}. Without an IOMMU, these are no-ops (the
kernel trusts the driver). With an IOMMU, they manipulate the translation table.

\section{Driver service layering}

Complex I/O stacks decompose into separate failure and authority domains:

\begin{verbatim}
Application
    -> socket endpoint
Network service
    -> packet endpoint
NIC driver
    -> MMIO / DMA / IRQ
Hardware
\end{verbatim}

The NIC driver knows about hardware registers and descriptor rings, the network
service knows about protocols and sockets, and the application knows about its
business logic. A crash in any layer is contained: the NIC driver crashes,
the network service sees \texttt{ServiceRestarted}, and reconnects to a fresh
instance.

\section{Why userspace drivers matter for critical software}

\begin{enumerate}
    \item \textbf{Driver bugs are contained.} A buffer overflow in a NIC driver
        corrupts that driver's address space, not the kernel's. The driver
        crashes (and restarts); the system stays up.

    \item \textbf{The privilege model is grant-by-grant.} A driver receives
        exactly the resources it needs. There is no ``driver has access to
        everything because it's in the kernel'' default.

    \item \textbf{Direct MMIO keeps the fast path fast.} Reading a device
        register is one load instruction. No syscall, no kernel transition,
        no context switch. The driver is in userspace, but the performance
        cost is near zero.

    \item \textbf{Restart is automated and defined.} The supervisor handles
        driver restart (device reset, grant re-issuance, generation bump).
        Driver authors do not write recovery code.

    \item \textbf{Hardware-isolated DMA is planned.} The IOMMU integration path
        is designed into the architecture from the start. When IOMMU support is
        enabled, DMA-capable drivers gain hardware-enforced memory isolation
        without API changes.
\end{enumerate}

\endinput
